\paragraph{Matched joint-native scope and scheduling controls.}
We replace the external restoration oracle with VPR's native connection router on copied trees, updating occupancy after every attachment. The released sink sets and all other retained branches are identical to the earlier transactions. Fixed-target R1 recovery restores four blocker sinks and reaches the same 5.854659868 ns CPD and exact setup vector as the capacity oracle. More strongly, native sequential repair of both targets and the same four blocker sinks, without preset compatible target paths, reaches the same CPD and setup vector. Both policies accept both candidates, each improving 246 setup components without regression relative to the frozen R1 source. Their resources are 40 SB pairs, 970 CB sources, and planar wirelength 218191. Thus neither the external capacity restoration oracle nor preset compatible paths are necessary for this R1 gain \emph{conditional on the supplied released scope}. The scope was supplied from the earlier automatically generated transaction; these controls do not yet demonstrate autonomous native scope discovery. Attribution should shift toward current-state target selection, affected-tree scope and additional legal repair opportunities.

For the Tf diagnostic state at 6.203926706 ns, fixed-target native recovery in net-ID order and in current-slack order both produce legal 6.096180893 ns candidates. Both improve 254 setup components and regress one, so the strict policy accepts and the endpoint-vector policy rejects. Their resources are 122 SB pairs and 959 CB sources; planar wirelength is 228423 and 228422 respectively. In the slack-order arm, constrained connections are restored before the eight unconstrained connections; this still does not prevent the output-victim loss. The four native transactions take 0.165, 0.171, 0.668 and 0.652 s internally. Their complete experimental orchestration takes 329.640 s, including separate candidate serialization, independent checks, fresh FULL/STA, preparation and logging. Individual independent audits/serialization take approximately 27.5--27.9 s, and replacement evaluation wrappers approximately 38.0--40.3 s. These development measurements motivate state reuse or in-process verification; they are not a formal cold-start routing speedup comparison.

\paragraph{A local bound changes the next intervention.}
An exact zero-RC fixed-retained-tree query for the directly connected \texttt{PRDATA[7]} victim, before \emph{any} displaced sink is restored, gives 1.681467215 ns versus its original 1.220990892 ns. Therefore the observed 0.460476324 ns physical connection loss cannot be removed solely by reordering the already released sinks while retaining the target paths and unchanged background. A further optimistic relaxation retaining only the fixed target trees gives 0.849339293 ns. This is a physical-connection reference, not a globally attainable CPD or STA bound. Its actual path conflicts identify four additional background nets and six sink branches. A feedback-derived transaction releases exactly those branches and restores the victim first. Four connections are recovered, but native recovery of net 1960/pin 1 continues to return a capacity-conflicting path through the tested pressure range; the whole transaction is rolled back. No candidate CPD is assigned to this failed transaction. The next discriminating test requires its partial-state witness and an exact capacity query, to distinguish insufficient released scope from an ineffective native soft congestion price. No global infeasibility or no-benefit conclusion follows from the guarded native failure.
